% Paper-ready chapter fragment; uses booktabs. Evidence cutoff: 2026-10-09 CST.
% Sources: demand53005_02/{demand_results,service_decomposition}.json,
% refresh53005_01/fixed_results.json, action_rollout/{warm_validation,three_states}.json.
% Separate GPU groups are not pooled. The bounded GPU probe is UNRUN.
\section{Does calibrated batch evolution provide incremental decision value?}
\label{sec:d-calibrated-decision}

We study whether state-dependent prefill budgets improve complete-service
outcomes beyond a tuned fixed budget, rather than merely advancing one request's
first output. The deployment is BF16 OLMoE on one RTX PRO 6000 Blackwell using
vLLM 0.26, a 90\% GPU-memory target, 192 sequence slots, and 588,032 usable KV
tokens. The seen development trace contains 256 open-loop requests, 538,597
prompt tokens, and 114,688 fixed output tokens. Only the prefill budget changes;
native admission, FCFS order, decode progress, recovery, precision, and kernels
remain common. Our research contract requires TTFT $\leq4$ s, per-request maximum
generation gap $\leq100$ ms, and completion latency $\leq20$ s. These thresholds
are research assumptions, not validated application SLOs. Goodput is the number
of jointly qualifying requests divided by the full elapsed time, including drain
and control costs.

\begin{table}[ht]
\centering\small
\begin{tabular}{p{0.25\linewidth}p{0.32\linewidth}p{0.34\linewidth}}
\toprule
Evidence and comparison & Local or waiting effect & Complete or cohort outcome \\
\midrule
GPU: demand vs. fixed1024; two contemporaneous pairs
& Mean TTFT $+1.637/+1.632$ s; generation span $-0.878/-0.951$ s
& Mean flow $+0.759/+0.681$ s; qualified $-32/-30$; goodput $-19.89/-18.72\%$ \\
GPU: fixed2048 vs. fixed1024; separate contemporaneous pairs
& TTFT p95 $-81.62/-83.08\%$; output throughput $+0.765/+0.837\%$
& Qualified $178/178$ vs. $183/183$; goodput $-1.99/-1.92\%$ \\
Model only: temporary2048 vs. 1024, high-decoder state
& First step $+5.951$ ms; head first output $-83.238$ ms
& Head completion $-16.702$ ms; known-cohort mean flow $+0.966$ ms; qualified $126\to126$ \\
Frozen bounded GPU probe
& Unrun: zero warmups and formal runs
& No measured intervention or service effect \\
\bottomrule
\end{tabular}
\caption{Evidence for the decision-value question. Slashes retain both
run-level comparisons, not confidence intervals. GPU groups have their own
contemporaneous controls and are not pooled. The model row drains 133 currently
known requests and excludes future arrivals; it is not full-trace goodput.}
\label{tab:d-decision-value}
\end{table}

Under the research SLO of 4 s TTFT, 100 ms maximum generation gap, and 20 s
completion, the demand rule reduced measured goodput by 19.89\% and 18.72\% in
two reversed GPU pairs. Despite 727 and 726 actual interventions, it reduced
mean generation span by 0.878/0.951 s but increased TTFT by 1.637/1.632 s,
raising mean flow time by 0.759/0.681 s. It gained 13 qualifying requests but
lost 45/43; all qualification failures concerned completion. In a separate
contemporaneous comparison, fixed2048 also reduced goodput by 1.99\%/1.92\%
relative to fixed1024; these groups are not pooled as controls. A simulator
calibrated from warmup reproduced the demand rule's adverse direction in a
seen-data explanatory check. Yet, at a state with 131 decoders, temporarily
using 2048 until the head's first output predicted only a local tradeoff: the
first step grew by 5.951 ms, head first output advanced by 83.238 ms, and head
completion advanced by 16.702 ms, while mean flow increased by 0.966 ms and
qualifications remained 126. This prediction drains only currently known
requests, excluding future arrivals. It does not establish full-service value
beyond fixed1024. The frozen GPU probe remains unrun. Fixed output lengths do
not establish identical content or quality. Two repeats per arm support
exploratory comparisons, not statistical confirmation.

All 256 requests completed in each measured GPU run, with no failures,
rejections, timeouts, unfinished requests, or preemptions. The demand rule
rescued the same 13 of the baseline's 73 completion violations but introduced
45/43 new completion violations; no TTFT or gap violation explains its loss.
Its total elapsed time increased by 1.479/1.412 s. Thus TTFT feasibility is
insufficient to protect completion slack. In the separate fixed-budget group,
2048 lost five qualifications and gained none despite draining faster.
These are complete-policy observations, not matched-state causal effects.

Calibration used 5,715 existing fixed-budget warmup steps, the original two
affine cost regimes, and an estimated 270.29~$\mu$s busy host-loop gap. The demand
policy retained its original predictor. Formal outcomes were not used to fit
the model. An unaligned compilation warning in the first warmup was retained,
so the fit is not claimed to be compilation-free. The calibrated model predicts
a $+0.620$ s mean-flow penalty for demand, versus $+0.759/+0.681$ s measured.
This explains a failure direction on already-seen data; it does not validate
counterfactual action ranking. All three preselected low/middle/high decoder
states prefer the larger action by conditional drain time, with unchanged
qualification counts. They therefore provide no demonstrated state-dependent
preference reversal or advantage that a simple larger budget cannot explain.

\paragraph{Decision and remaining limit.}
We retain fixed1024 and stop developing the current dynamic-budget route as
an independent method contribution. The supported material is an adaptation
of known scheduling principles, ordinary cost calibration, and a scoped
negative result about waiting--generation tradeoffs. It is not a proof that
state-aware chunking is generally ineffective, nor a sufficient independent
measurement contribution. The remaining bounded action probe was not executed
because the shared GPU was occupied; storage had recovered, no D process was
submitted, and no lock wait was started. Its frozen two-warmup, F/X/X/F design
is preserved. A positive future probe would establish only action value versus
its contemporaneous fixed1024 control; comparison with a strong fixed-large
policy under the same input and timing protocol would still be necessary to
establish incremental state value. Historical fixed2048 results cannot fill
that gap. No online rollout overhead or deployment benefit has been measured.
